% Generated by tools/textbook/export_listings.js from tools/concepts/assessments.json: do not edit.
\begin{chapstatus}{Il filo del capitolo}
Ogni rettangolo disegnato dalla Lezione 02 è finito in un frame buffer che un controllore invia allo schermo molte volte al secondo. Prossimo passo: gli altri dispositivi: display, tastiere, stampanti.
\par\medskip\textbf{Biglietto d'uscita} (da due a quattro minuti, da soli, su carta):
\begin{enumerate}[nosep,leftmargin=2em]
\item Perché un dispositivo non è mai collegato direttamente al processore? Dite che cosa fa il suo controllore e come il processore dialoga con esso.
\item Uno schermo di 320 x 200 pixel usa 16 bit per pixel ed è ridisegnato 30 volte al secondo. Quanta memoria serve per il frame buffer, e quanti byte al secondo devono raggiungere lo schermo?
\end{enumerate}
\end{chapstatus}
